\section{A complete 663-bit equational consistency packet}
\label{sec:e667}

The alternative-formalism branch yields a complete \textbf{663-bit}
consistency packet with seven equations, nine operation symbols,
150 equation-term nodes, and the native-coordinate query $Q(d)=d$.
It is obtained in two independently justified steps.  First, combining
the 236-state archive's derived-truth argument with the separate
$\mathcal E_{692}$ packet's structural witnesses, Boolean multiplexer,
and unary selector gives the 667-bit intermediate $\mathcal E_{667}$.
Second, replacing the four-node Boolean collapse query $N(d,d)=d$
by the three-node native collapse query $Q(d)=d$ saves four more bits.
The latter query attains the minimum possible endpoint length among
distinct ground endpoints in this unchanged signature and codec.
Both reductions use the existing codec; neither charges a new opcode
or a definition supplied outside the packet.

The first reduction has a particularly strong relative conclusion:
\begin{equation}
 \operatorname{Mod}(\mathcal E_{667})
   =\operatorname{Mod}(\mathcal E_{692}).
 \label{eq:e667-models}
\end{equation}
Thus it preserves every equational consequence, in the identical
signature, not merely the designated consistency question.

The first contribution is the explicit combination and its proof for the
\emph{changed} multiplexer closure.  The prior $\mathcal E_{688}$ argument
proved truth redundant for a different closure expression and a larger
signature; its conclusion did not automatically apply to
$\mathcal E_{692}$.  We check the relevant Boolean cofactor again.
The new query reduction then uses only native Head/Tail to establish
carrier collapse.  The inherited 692-bit packet itself and its model
construction retain their attribution.  No binary Turing-machine bound
follows from the 663 bits: an interpreter, complete search, and blank
initialization must also be charged.  We retain the complete E667
intermediate to separate the two arguments and their exact accounting.

\subsection{The literal theory and its designated query}

The ordered signature is
\[
 d^0,N^2,V^2,P^2,A^2,Q^1,W^2,p_E^2,p_\varepsilon^2.
\]
The first six equations, universally quantified in all displayed
variables, are
\begin{align}
 N(N(x,N(N(y,x),x)),N(y,N(z,x)))&=y,
       \label{eq:e667-sh1}\\
 V(P(x,s),d)&=x,\label{eq:e667-head}\\
 V(P(x,s),Q(j))&=V(s,j),\label{eq:e667-tail}\\
 p_\varepsilon(s,P(a,b))
     &=N(p_\varepsilon(s,a),p_\varepsilon(s,b)),
       \label{eq:e667-nand}\\
 p_\varepsilon(s,A(i,j))
     &=p_E(V(s,i),V(s,j)),\label{eq:e667-atom}\\
 p_\varepsilon(s,Q(a))
     &=N\bigl(p_\varepsilon(P(W(s,a),s),a),
               p_\varepsilon(P(z,s),a)\bigr).
       \label{eq:e667-binder}
\end{align}
Equation~\eqref{eq:e667-sh1} is the published Sheffer single axiom Sh1.
Theorem 2 of McCune, Veroff, Fitelson, Harris, Feist, and Wos,
\emph{Short Single Axioms for Boolean Algebra}, supplies its Boolean
algebra interpretation and a 66-step equational derivation of the
Sheffer basis.\footnote{Primary paper:
\url{https://www.cs.unm.edu/~mccune/papers/basax/v12.pdf};
the authors' supplementary material is at
\url{https://www.cs.unm.edu/~veroff/BA/}.}
For the Boolean algebra determined by this reduct, $N$ is NAND;
write its derived top and bottom as $1$ and $0$.
\emph{We do not assume $d=1$ at this stage.}

The seventh equation uses these displayed abbreviations, all expanded
in the delivered AST and bit packet:
\begin{align*}
 h&=Q(b),& n_z&=P(b,z),&m_z&=A(b,z),\\
 R_b(x,y)&=p_\varepsilon(P(y,P(x,b)),V(b,d)),&&&
\end{align*}
\begin{align*}
 r&=R_b(h,n_z),& e&=p_E(m_z,n_z),&f&=p_E(m_z,z),\\
 a&=p_E(W(x,d),x),& b_0&=p_E(x,h),&c&=R_b(x,y),\\
 u&=p_E(y,h),&v&=R_b(P(b,y),y).&&
\end{align*}
Here $b$ is the descriptor variable, while $b_0$ names a Boolean
predicate value.  Define the Boolean terms
\begin{align*}
 K&=N(c,N(a,N(d,b_0))),\\
 H&=N\bigl(N(N(d,u),K),N(u,v)\bigr),\\
 B^-&=N\bigl(r,N(d,N(e,N(d,f)))\bigr),\\
 F(d,r,e,f,a,b_0,c,u,v)&=N(B^-,N(d,H)).
\end{align*}
The seventh equation is exactly
\begin{equation}
             d=F(d,r,e,f,a,b_0,c,u,v).
 \label{eq:e667-closure}
\end{equation}
The designated ground equality whose derivability is searched for is
\begin{equation}
                         N(d,d)=d.
 \label{eq:e667-query}
\end{equation}
The inherited $\mathcal E_{692}$ consists of these seven equations
together with
\begin{equation}
                         d=N(d,N(d,d)).
 \label{eq:e667-deleted}
\end{equation}
This additional equation has six term nodes and states $d=1$ once
Sh1 is available.

\subsection{The truth-forcing argument survives the multiplexer}

\paragraph{Boolean lemma.}
In every Boolean algebra, treating all nine arguments independently,
\begin{equation}
 \neg d\wedge F(d,r,e,f,a,b_0,c,u,v)=\neg d\wedge r.
 \label{eq:e667-cofactor}
\end{equation}
In particular, $r=1$ and $d=F$ imply $d=1$.

\paragraph{Proof.}
At formal argument $d=0$, every subterm $N(d,t)$ is $1$.
Consequently $B^-=N(r,1)=\neg r$ and
$F=N(\neg r,1)=r$, regardless of the internal value of $H$.
For an arbitrary Boolean algebra, project evaluated Boolean values
to the principal interval $[0,\neg d]$ by $t\mapsto t\wedge\neg d$.
This is a Boolean homomorphism when the interval uses its relative
complement and top.  It sends $d$ to zero, so the preceding
calculation proves~\eqref{eq:e667-cofactor}.
If $r=1$ and $d=F$, its left side is zero and its right side is
$\neg d$, proving $d=1$.
Only the already evaluated Boolean values are projected;
there is no assertion that $P,V,W$ commute with this projection.
\hfill$\square$

An alternative short argument for the only required case $r=1$
is that $B^-=\neg N(d,X)=d\wedge X$ for
$X=N(e,N(d,f))$.  Hence $B^-\leq d$, and
$F=\neg(B^-\wedge N(d,H))\geq\neg d$.  The equality $d=F$
then implies $\neg d\leq d$, which forces $d=1$.

\paragraph{A concrete ground instance supplies $r=1$.}
Let
\[
                  k=P(d,P(d,d)),\qquad b_*=P(k,d).
\]
The NAND-class equation gives, for every native argument $s$,
\[
 p_\varepsilon(s,k)
 =N\bigl(p_\varepsilon(s,d),
           N(p_\varepsilon(s,d),p_\varepsilon(s,d))\bigr)=1.
\]
The last equality is $t\mathbin{\mathrm{NAND}}\neg t=1$.
It imposes no assumption on the class denoted by $d$.
By the native head equation, $V(b_*,d)=k$.  In the closure
equation substitute $b=b_*$ and $x=y=z=d$.  Its first predicate is
\[
 r=p_\varepsilon(P(P(b_*,d),P(Q(b_*),b_*)),k)=1.
\]
Applying the Boolean lemma proves $d=1$ and therefore the omitted
equation~\eqref{eq:e667-deleted}.

\paragraph{Same-model theorem.}
Every $\mathcal E_{667}$-model is an $\mathcal E_{692}$-model,
and conversely.  The forward inclusion is the derivation just
given; the reverse inclusion deletes an axiom.  Equational
soundness and completeness yield identical equational consequences,
including identical derivability status for~\eqref{eq:e667-query}.
This proves~\eqref{eq:e667-models}.  The new argument is ordinary
mathematics formalizable in ZFC.  The artifact contains its explicit
ground substitution and its Boolean identity, not a newly emitted
rewrite derivation expanding every use of the published Sh1 theorem.

\subsection{Consistency calibration and the scope of the inherited proof}

The inherited $\mathcal E_{692}$ calibration is an ordinary model
construction over the equality-free empty-input closure schema
(raw EPSC).  Its full proof is retained as
\texttt{reference/e692\_compatibility.tex}.  The crucial points are
as follows.

After a homomorphism from the nontrivial Boolean reduct to two
truth values, define $E(x,y)$ by the resulting value of $p_E(x,y)$
and class incidence by that of $p_\varepsilon$.
The native equations retain literal equality on the original carrier.
The binder equation implies
\[
 \varepsilon(s,Q(a))\quad\Longleftrightarrow\quad
        \exists z\,\neg\varepsilon(P(z,s),a),
\]
with the selected $W(s,a)$ realizing a false instance whenever
one exists.  To see this, write the selected truth value as $w$
and the arbitrary one as $u_z$.  The equations
$q=\neg(w\wedge u_z)$ for all $z$, including the selected value,
give $q=\neg w$ and $w\leq u_z$ for every $z$.

Selectors $\nu_0=d$, $\nu_{i+1}=Q(\nu_i)$ and the two native
stack equations recover every finite coordinate.  The class constructor
$A$ represents atomic membership, $P$ represents NAND, and
$Q(P(a,a))$ represents ordinary existential quantification.
Thus every finite pure-membership formula with finitely many parameters
has a descriptor $b$.  The head code is built before that descriptor,
with the two relation variables placed before it on the stack.

Once $d=1$ is established, the closure equation is exactly
\[
 (r\wedge e\wedge\neg f)
 \ \lor\
 \bigl((a\lor\neg c\lor u)
       \wedge(\neg b_0\lor\neg c\lor u)
       \wedge(\neg u\lor v)\bigr).
\]
For a locally subset-bounded represented relation, a bound at $h=Q(b)$
makes one bad-branch instance false.  All hull instances must then hold.
They give closure at empty inputs and at members of $h$, together
with unrestricted relation support for every member of $h$.
These are exactly raw EPSC.

Conversely, the inherited proof starts with a countable raw-EPSC
structure and takes a free tagged carrier with constructors $P,A,Q$.
The integer weights
\[
 w(Q(b))=w(b)+1,\quad
 w(P(b,z))=w(b)+w(z)+2,\quad
 w(A(b,z))=w(b)+w(z)+3
\]
order raw-value choices: a hull or unbounded input at $Q(b)$,
a relation witness at $P(b,z)$, then a separating member at $A(b,z)$.
Descriptor parameters descend to proper subtrees, so the choices
are well founded.  The independent class and selector projections
respect their own syntactic recursions.  The class $d$ is false,
leaving $W(x,d)$ free to select a member of every nonempty raw object.
Transporting a countably infinite Boolean algebra to the carrier,
with its top at $d$, supplies the separating equational model.

The raw-EPSC/ZFC comparison is the expressly inherited ordinary
calibration from ProveIt's Sources 14 and 08; Source 14 proves the
EPSC/PSC comparison and records the hereditary interpretation and
Foundation/Choice dependencies.\footnote{Pinned source dossier:
\url{https://github.com/VladimirReshetnikov/ProveIt/blob/ce192939372ac7944fe2a0208c68a96d49fb5e3b/SetTheory/Cardinals/docs/reports/set-theory/zfc-closure-kernels/intake-146-source-14.md}.
A retrieved snapshot is retained in \texttt{reference/raw\_epsc\_source14.md}.}
Composing that calibration with the new same-model theorem yields
\begin{equation}
 \mathrm{ZFC}\vdash
 \left(\mathcal E_{667}\nvdash N(d,d)=d
          \ \longleftrightarrow\ \operatorname{Con}(\mathrm{ZFC})\right).
 \label{eq:e667-calibration}
\end{equation}
The present work does not claim a new formal-kernel verification of
the entire inherited interpretation, a PA proof-code transformer,
or a binary TM implementing this packet.

\subsection{Two tempting signature identifications fail}

The successful identification of structural constructors with raw
closure witnesses should not be extrapolated without a compatibility
proof.  Two further identifications have direct obstructions.

\paragraph{Identifying $W$ with $V$ forces a common member.}
Assume a separating model of the theory obtained by this identification
and pass to its two-valued raw-EPSC interpretation.  Put
\[
 i=Q(d),\qquad j=Q(Q(d)),\qquad a=Q(A(i,j)).
\]
For arbitrary $x,z,s$, one binder expansion and the stack identities give
\begin{align*}
 \varepsilon(P(z,P(x,s)),a)
 &\Longleftrightarrow
   \exists w\,\neg E(z,x)\\
 &\Longleftrightarrow\neg E(z,x).
\end{align*}
Nonemptiness of the carrier justifies the second equivalence.
Therefore the outer complemented binder, when applied to environment
$P(x,s)$ and class $a$, must choose a member of $x$ whenever $x$
is nonempty.  But $W=V$ and the native tail equation give
\[
 W(P(x,s),a)=V(P(x,s),Q(A(i,j)))=V(s,A(i,j)).
\]
For fixed $s$ the selected object is independent of $x$, so it must
belong to \emph{every} nonempty raw object.

Raw EPSC has disjoint nonempty objects: its established bootstrap
gives an empty object $e$, a nonempty power object of $e$, and
full-profile singleton objects for each.  The first singleton contains
only empty-profile objects; the second contains only nonempty-profile
objects.  Each contains its defining object, so both are nonempty
and they are disjoint.  No common member exists.  This contradiction
shows that the $W=V$ specialization has no separating model and its
designated equation is derivable.  The obstruction arises from the
native tail equation, not from an unsuccessful finite search.

\paragraph{Identifying $p_E$ with $p_\varepsilon$ creates a universal set.}
The ground class $k=P(d,P(d,d))$ used above has
$p_\varepsilon(s,k)=1$ for every $s$.  If the two predicate operations
are identical, the raw relation therefore satisfies $E(s,k)$ for
every $s$.  Such a universal object is incompatible with raw EPSC
by a direct two-instance argument.

First the constantly false relation has a locally bounded output
class, and EPSC support produces an empty object $e$.
Next use the relation $R(x,y)\Longleftrightarrow\neg E(y,y)$.
Its outputs are locally subset-bounded by the universal object $k$.
EPSC closure at $e$ gives
$\neg E(y,y)\Rightarrow E(y,h)$; support gives the converse.
Substitution $y=h$ yields $E(h,h)\Longleftrightarrow\neg E(h,h)$.
Thus this identification also destroys the separating models.
These arguments explain why the remaining native witness and predicate
symbols cannot simply be merged with their nearest apparent analogues.

\subsection{Exact accounting, verification, and next implementation boundary}

The seven equation-node counts are
\[
                   14,6,9,12,12,17,80,
\]
whose sum is 150.  Nine operations and four equation-local variables
still require four-bit term opcodes.  The complete count is
\[
 \underbrace{8+3\cdot9+8}_{43\text{ header}}
 +\underbrace{7+1+4\cdot150}_{608\text{ equations}}
 +\underbrace{4\cdot4}_{16\text{ endpoints}}
 =\boxed{667\text{ bits}}.
\]
The exact bit-character SHA-256, excluding the final newline, is
\begin{center}\scriptsize\ttfamily
1075854399e1c3428e1c178b51bac47e428b8ef010b59e0312a3c06d409e55b8.
\end{center}
Packed transport occupies 84 bytes, with five trailing zero padding
bits.  Those padding bits are not part of the logical packet.

\begin{center}
\begin{tabular}{lrrrr}
\hline
Packet & Operations & Equations & Equation nodes & Logical bits\\
\hline
Inherited $\mathcal E_{688}$ & 12 & 7 & 153 & 688\\
Inherited $\mathcal E_{692}$ & 9 & 8 & 156 & 692\\
Combined $\mathcal E_{667}$ & 9 & 7 & 150 & 667\\
Native-query $\mathcal E_{663}$ & 9 & 7 & 150 & 663\\
\hline
\end{tabular}
\end{center}

Run \texttt{python3 -B build\_e667.py} in this companion directory.
The independent builder first reconstructs the exact inherited
$\mathcal E_{688}$ and $\mathcal E_{692}$ packets and checks their
literal supplied bits and pinned hashes.  It then verifies that
$\mathcal E_{667}$ deletes only the indicated truth equation,
round-trips the full new packet, rejects all 667 strict truncations
and three malformed controls, and checks all 512 Boolean assignments
for the cofactor and truth-forcing implication.  A further 256
assignments check the closure translation after $d=1$.
The finite checks establish transcription and Boolean identities;
the proof above handles arbitrary Boolean algebras and unbounded models.

The complete derivability target admits fair enumeration of ground
contextual rewrites in either orientation.  A certificate must supply
ground values for the union of variables on both sides of each used
equation, including variables absent from the side being matched.
Taking the least congruence of all such ground instances gives the
usual countable separating ground-term algebra whenever the two
endpoints are not connected.  The constant $d$ ensures that the
ground-term domain is nonempty.

The unresolved implementation cost is substantial.  The repository's
separate equational-generator experiment reports a 299-working-state
control, before a literal ZFC query and its writer, with a 171-state
isolated emitter; those are a different input basis and are not a
compilation of $\mathcal E_{667}$.\footnote{Pinned cost report:
\url{https://github.com/VladimirReshetnikov/ProveIt/blob/13ec39ed43148b1807f51ed6e6528e1c84659562/Computability/BusyBeaver/Research/ZFC/equational_generator_cost.md}.}
A competitive next construction therefore needs a smaller interpreter
architecture in addition to the compact axioms.  E667 is an exact,
calibrated input for that work; the main machine construction in this
report has its own independently counted finite table.

\subsection{A further reduction to 663 bits: query native noncollapse}
\label{sec:e663}

The truth-value query $N(d,d)=d$ is not the shortest available ground
query with the required calibration.  Keep all seven equations of
$\mathcal E_{667}$ and the complete ordered signature unchanged, but
replace the designated query by
\begin{equation}
                            Q(d)=d.
 \label{eq:e663-query}
\end{equation}
The resulting complete consistency packet has \textbf{663 bits}.
This reduction uses the native coordinate equations, independently of
the Boolean truth-forcing optimization.

\paragraph{Query-equivalence theorem.}
For every algebra satisfying the seven equations, each of
$Q(d)=d$ and $N(d,d)=d$ holds if and only if its carrier is a
singleton.  Consequently the two queries have the same equational
derivability status from the seven equations.

\paragraph{Proof.}
Suppose $Q(d)=d$.  The native head and tail equations, with the
tail index $j=d$, give for every $x,s$
\[
 x=V(P(x,s),d)=V(P(x,s),Q(d))=V(s,d).
\]
Fixing $s=d$ shows that every element $x$ equals $V(d,d)$, so the
carrier is a singleton.  No Boolean identity or semantic quotient
is used in this implication.

Sh1 makes $N$ NAND in a Boolean algebra on this same carrier.
If $N(d,d)=d$, then $\neg d=d$, which implies
$0=d\wedge\neg d=d$ and $1=d\vee\neg d=d$.
Thus $0=1$ and the carrier is again a singleton.  Conversely both
equalities hold in every singleton algebra.  The two equalities
therefore hold in exactly the same models of the seven equations.
Equational soundness and completeness give equivalence of their
derivability.  Equivalently, use the countable ground-term congruence
model to witness either nonderivability; the constant $d$ guarantees
that this model is nonempty.
\hfill$\square$

The E667 consistency theorem therefore gives, with its same inherited
ordinary raw-EPSC/ZFC calibration,
\[
 \mathrm{ZFC}\vdash
 \left(\mathcal E_{667}\nvdash Q(d)=d
              \ \longleftrightarrow\ \operatorname{Con}(\mathrm{ZFC})\right).
\]
The label $\mathcal E_{663}$ denotes this new \emph{packet}; the
seven-equation axiom set itself is unchanged.  This distinction matters
because the four-bit saving is in the encoded goal, not the axioms.

The header still has 43 bits and the equations have 608 bits.  The new
endpoints have three term nodes, so they require 12 rather than 16 bits:
\[
                       43+608+12=\boxed{663\text{ bits}}.
\]
Packed transport occupies 83 bytes, with one trailing zero padding bit.
The exact bit-character SHA-256, excluding the final newline, is
\begin{center}\scriptsize\ttfamily
6c3b72c39c81aa5bbf028e654c8056e7e86376cdeecdd56943e5c3e23de1a215.
\end{center}
Run \texttt{python3 -B e663/build\_e663.py} from the E667 companion
directory.  It reconstructs the pinned E667 packet, verifies that the
header and all axiom bits are literally unchanged, and emits
\texttt{E663.bits}, \texttt{E663.bin}, \texttt{E663.bitlength},
\texttt{E663.json}, and an exact verification receipt.
All 663 strict truncations and three malformed controls are rejected
by the independent decoder.

\paragraph{A limited endpoint optimum.}
The signature has exactly one constant, $d$.  Thus its only one-node
ground term is $d$.  Two syntactically distinct ground endpoints must
have at least three nodes in total: one endpoint has at least two
nodes, and the other has at least one.  The pair $Q(d),d$ attains this
bound.  Since every node costs four bits in the unchanged codec,
\textbf{12 endpoint bits are optimal among distinct ground endpoint
pairs in this fixed signature and codec}.  This is not a claim that
663 total bits are globally minimal: the axioms, signature, codec,
or interpreter could be redesigned.  No machine-state bound is
inferred from the shorter query.
